\section*{Chapitre \ref{ch:b2:reaction-free-energy} --- Entropie et enthalpie libre de réaction}

\begin{solution}{exo:b2:reaction-free-energy:1}
(a) Trois moles de gaz deviennent liquides~: $\Delta_r S^\circ < 0$. (b) Un gaz est
formé~: $> 0$. (c) Une molécule de gaz en donne deux~: $> 0$. (d) Des ions en solution
forment un cristal~: $< 0$. (e) Fusion~: $> 0$. (f) $\Delta\nu_{\text{gaz}} = 0$~:
petite, et ici positive (\qty{+20}{J/(K.mol)} d'après les tables, \ce{HCl}
étant moins symétrique que \ce{H2} et \ce{Cl2}).
% ledger: s0:HCl_g, s0:H2_g, s0:Cl2_g
\end{solution}

\begin{solution}{exo:b2:reaction-free-energy:2}
$\Delta_r S^\circ = 2(192.77) - 191.61 - 3(130.68) = \qty{-198.1}{J/(K.mol)}$~;
$\Delta_r G^\circ = -91.9 - 298.15 \times (-0.1981) = \qty{-32.8}{kJ/mol}$~:
la réaction est \omterm{def:b2:reaction-free-energy:exergonic}{exergonique} à \qty{298}{K}, bien que le terme entropique (\qty{+59.1}{kJ/mol})
lui soit défavorable.
\end{solution}

\begin{solution}{exo:b2:reaction-free-energy:3}
$\Delta_r S^\circ = 69.95 - 130.68 - \tfrac12(205.15) = \qty{-163.3}{J/(K.mol)}$~;
$\Delta_r G^\circ = -285.83 + 298.15 \times 0.1633 = \qty{-237.14}{kJ/mol}$.
C'est la réaction de formation de l'eau liquide (une mole à partir des éléments dans
leur état de référence)~: son \omterm{def:b2:reaction-free-energy:reaction-gibbs}{enthalpie libre standard de réaction} est donc par
définition $\Delta_f G^\circ(\ce{H2O}, \text{l})$~; les tables donnent la même
valeur.
\end{solution}

\begin{solution}{exo:b2:reaction-free-energy:4}
Environ \qty{42}{J/(K.mol)} à \qty{298}{K} et \qty{87}{J/(K.mol)} à
\qty{1000}{K} (liquide). Sauts~: \qty{10.6}{J/(K.mol)} à \qty{692.7}{K} et
\qty{97.7}{J/(K.mol)} à \qty{1180.2}{K}. $\Delta_{\text{fus}}H = 10.6 \times
692.7 = \qty{7.3}{kJ/mol}$ et $\Delta_{\text{vap}}H = 97.7 \times 1180.2 =
\qty{115}{kJ/mol}$.
\end{solution}

\begin{solution}{exo:b2:reaction-free-energy:5}
$\Delta_r H^\circ = \qty{44.00}{kJ/mol}$, $\Delta_r S^\circ = 188.84 - 69.95 =
\qty{118.9}{J/(K.mol)}$, $\Delta_r G^\circ = 44.00 - 298.15 \times 0.1189 =
\qty{+8.55}{kJ/mol}$. À \qty{298}{K}, la transformation n'est pas réversible (l'eau
liquide et sa vapeur sous \qty{1}{bar} ne sont pas en équilibre), donc $\Delta S \ne
Q/T$~: $\Delta_r H^\circ/T = \qty{147.6}{J/(K.mol)}$. La valeur positive de
$\Delta_r G^\circ$ signifie qu'à \qty{298}{K} l'eau ne s'évapore pas dans
une atmosphère de vapeur pure sous \qty{1}{bar}~: la vapeur se condense. Les deux
sont en équilibre là où $\Delta_r G^\circ = 0$~: $T = 44\,004/118.9 =
\qty{370}{K}$, à \qty{3}{K} près de la température d'ébullition normale
(\qty{373}{K})~: l'\omterm{def:b2:reaction-free-energy:ellingham-approximation}{approximation d'Ellingham} est bonne sur \qty{75}{K}.
\end{solution}

\begin{solution}{exo:b2:reaction-free-energy:6}
$T_i = 57\,100/175.7 = \qty{325}{K}$. Au-dessous, $\Delta_r G^\circ > 0$ et
\ce{N2O4}, incolore, est favorisé~; refroidir le tube déplace le gaz vers
\ce{N2O4}, et la couleur brune de \ce{NO2} s'atténue (le lien quantitatif avec
la composition est établi au \cref{ch:b2:equilibrium-shifts}).
\end{solution}

\begin{solution}{exo:b2:reaction-free-energy:7}
$\Delta_r S^\circ = -\dd\Delta_r G^\circ/\dd T = -(-4.0 + 12.0)/100 =
\qty{-0.080}{kJ/(K.mol)} = \qty{-80}{J/(K.mol)}$ et $\Delta_r H^\circ =
\Delta_r G^\circ + T\Delta_r S^\circ = -12.0 + 300 \times (-0.080) =
\qty{-36.0}{kJ/mol}$. Vérification~: $\Delta_r G^\circ/T$ passe de $-0.0400$ à
$\qty{-0.0100}{kJ/(K.mol)}$, pente $3.00 \times 10^{-4}$ par kelvin, égale à
$-\Delta_r H^\circ/T^2 = 36.0/346^2 = 3.00 \times 10^{-4}$ à la température
moyenne géométrique $\sqrt{300 \times 400} = \qty{346}{K}$.
\end{solution}

\begin{solution}{exo:b2:reaction-free-energy:8}
$\Delta_r S^\circ = 186.25 - 5.74 - 2(130.68) = \qty{-80.8}{J/(K.mol)}$,
$\Delta_r G^\circ = -74.87 + 298.15 \times 0.0808 = \qty{-50.8}{kJ/mol}$,
c'est le $\Delta_f G^\circ$ tabulé du méthane. Les deux termes étant de même
signe, le méthane devient instable au-dessus de $T_i = 74\,870/80.8 \approx
\qty{926}{K}$ (\omterm{def:b2:reaction-free-energy:ellingham-approximation}{approximation d'Ellingham})~: le méthane chaud se craque en carbone et
dihydrogène, une voie d'accès à l'un et à l'autre.
\end{solution}

\begin{solution}{exo:b2:reaction-free-energy:9}
Substitution~: \ce{CH3CHBrCH3 + OH- -> CH3CH(OH)CH3 + Br-}~; élimination~:
\ce{CH3CHBrCH3 + OH- -> CH2=CHCH3 + H2O + Br-}. La différence des deux
\omterm{def:b2:reaction-free-energy:reaction-gibbs}{enthalpies libres standard de réaction} vaut $\Delta(\Delta_r G^\circ) = 40 - T \times
0.090$, négative au-dessus de $40\,000/90 \approx \qty{440}{K}$. L'élimination produit
une molécule de plus que la substitution~: son entropie est plus grande, et le terme
entropique croît proportionnellement à $T$.
\end{solution}

\begin{solution}{exo:b2:reaction-free-energy:10}
Chacune des $N_A$ molécules a 2 orientations~: $W = 2^{N_A}$, donc $S =
k_B\ln 2^{N_A} = R\ln 2 = \qty{5.76}{J/(K.mol)}$. Le cristal n'est pas
parfaitement ordonné~: à basse température, les molécules sont figées dans leurs
orientations aléatoires et ne peuvent atteindre l'unique configuration ordonnée que suppose
le troisième principe.
\end{solution}

\begin{solution}{exo:b2:reaction-free-energy:11}
Intégrer $\dd\Delta_r S^\circ/\dd T = c/T$~: $\Delta_r S^\circ(T) = \Delta_r
S^\circ(T_0) + c\ln(T/T_0)$~; la \omterm{thm:b2:reaction-enthalpy:kirchhoff}{loi de Kirchhoff} donne $\Delta_r H^\circ(T) = \Delta_r
H^\circ(T_0) + c(T - T_0)$~; retrancher $T\Delta_r S^\circ(T)$. Pour le calcaire à
\qty{1100}{K}~: Ellingham donne $178.32 - 1100 \times 0.16064 = \qty{1.62}{kJ/mol}$~;
correction $c[T - T_0 - T\ln(T/T_0)] = -0.00195 \times (801.85 - 1436.3) =
\qty{+1.24}{kJ/mol}$~: $\Delta_r G^\circ = \qty{2.86}{kJ/mol}$. L'erreur,
\qty{1.2}{kJ/mol}, vaut moins d'un pour cent de $\Delta_r H^\circ$, mais près de
la \omterm{def:b2:reaction-free-energy:inversion-temperature}{température d'inversion} elle décide du signe.
\end{solution}

\begin{solution}{exo:b2:reaction-free-energy:12}
(a) $\Delta_r H^\circ = -763.2 + 944.0 = \qty{+180.8}{kJ/mol}$, $\Delta_r
S^\circ = 353.2 + 205.15 - 50.62 - 2(223.08) = \qty{+61.6}{J/(K.mol)}$~;
$\Delta_r G^\circ(1000) = 180.8 - 61.6 = \qty{+119.2}{kJ/mol}$~: réaction \omterm{def:b2:reaction-free-energy:exergonic}{endergonique},
pas de chloration. (b) Pour \ce{C + O2 -> CO2}, $\Delta_r G^\circ(1000) =
-393.5 - 2.9 = \qty{-396.4}{kJ/mol}$. La somme, \ce{TiO2(s) + 2Cl2(g) + C(s)
-> TiCl4(g) + CO2(g)}, a pour $\Delta_r G^\circ = \qty{-277.2}{kJ/mol}$~: le carbone
consomme le dioxygène produit par la première réaction, et la réaction couplée est
fortement \omterm{def:b2:reaction-free-energy:exergonic}{exergonique}.
\end{solution}

\begin{solution}{pb:b2:reaction-free-energy:1}
\textbf{1.} \ce{CaCO3(s) -> CaO(s) + CO2(g)}.
\textbf{2.} $\Delta_r H^\circ = -635.09 - 393.51 + 1206.92 =
\qty{+178.3}{kJ/mol}$~: réaction \omterm{def:b2:reaction-enthalpy:exothermic}{endothermique}.
\textbf{3.} $\Delta_r S^\circ = 39.75 + 213.79 - 92.9 = \qty{+160.6}{J/(K.mol)}$,
positive car un gaz se forme à partir d'un solide.
\textbf{4.} $\Delta_r G^\circ = 178.32 - 298.15 \times 0.16064 =
\qty{+130.4}{kJ/mol}$.
\textbf{5.} Oui~: la décomposition exige $\Delta_r G < 0$, et ici l'\omterm{def:b2:reaction-free-energy:reaction-gibbs}{enthalpie libre de réaction}
est fortement positive~; le dioxyde de carbone de l'air, bien au-dessous de \qty{1}{bar},
abaisse $\Delta_r G$, mais de bien moins que \qty{130}{kJ/mol} (chapitre suivant).
\textbf{6.} $\Delta_r C^\circ_p = 42.80 + 37.13 - 81.88 =
\qty{-1.95}{J/(K.mol)}$, minuscule~: $\Delta_r H^\circ$ et $\Delta_r S^\circ$
varient à peine avec $T$.
\textbf{7.} $\Delta_r G^\circ(T) = 178.32 - 0.16064\,T$ (\unit{kJ/mol}).
\textbf{8.} \qty{+17.7}{kJ/mol} à \qty{1000}{K}~; \qty{-30.5}{kJ/mol} à
\qty{1300}{K}.
\textbf{9.} $T_i = 178\,320/160.64 = \qty{1110}{K}$.
\textbf{10.} $\Delta_r G^\circ/T = 178.32/T - 0.16064$, dont la dérivée vaut
$-178.32/T^2 = -\Delta_r H^\circ/T^2$.
\textbf{11.} En $T_i$, $c[T_i - T_0 - T_i\ln(T_i/T_0)] = -0.00195 \times
(811.9 - 1459.4) = \qty{+1.26}{kJ/mol}$~; $\Delta_r G^\circ$ s'annule
$1260/160.6 \approx \qty{8}{K}$ plus haut, vers \qty{1118}{K}~: moins
d'un pour cent.
\textbf{12.} À \qty{1000}{K}, $\Delta_r G > 0$~: la décomposition ne peut pas
avoir lieu (la chaux vive absorberait du dioxyde de carbone). À \qty{1300}{K}, $\Delta_r G <
0$~: le calcaire se décompose.
\textbf{13.} $\Delta_r G$ devient positive~: la chaux vive fixe le dioxyde de carbone
et redevient carbonate.
\textbf{14.} $\delta S_{\text{c}} = -\Delta_r G\,\dd\xi/T = 30\,506/1300 =
\qty{23.5}{J/K}$ par mole.
\textbf{15.} Oui~: le \cref{ch:b2:chemical-potential} montrera que
$\Delta_r G = \Delta_r G^\circ + RT\ln(p_{\ce{CO2}}/p^\circ)$, qui diminue
lorsque la pression du dioxyde de carbone est inférieure à \qty{1}{bar}~; la
décomposition commence alors au-dessous de $T_i$.
\textbf{16.} $n = 10^6/56.08 = \qty{1.783e4}{mol}$.
\textbf{17.} $1.783 \times 10^4 \times 178.32 = \qty{3.18e6}{kJ}$, soit
\qty{3.18}{GJ} par tonne.
\textbf{18.} $1.783 \times 10^4 \times 44.01 = \qty{785}{kg}$ de \ce{CO2}.
\textbf{19.} Chaleur à fournir $3.18 \times 10^6/0.50 = \qty{6.36e6}{kJ}$, donc
$6.36 \times 10^6/802.3 = \qty{7.93e3}{mol}$ de méthane, \qty{127}{kg}.
\textbf{20.} $7.93 \times 10^3 \times 44.01 = \qty{349}{kg}$ de \ce{CO2}.
\textbf{21.} $785 + 349 = \qty{1134}{kg}$ par tonne de chaux, dont 69\,\%
proviennent du calcaire.
\textbf{22.} Le dioxyde de carbone du calcaire est fixé par la stœchiométrie,
quelle que soit la façon dont on chauffe le four~; seule la part du combustible peut être réduite.
\textbf{23.} Le calcaire se décompose sous \qty{1}{bar} de dioxyde de carbone au-dessus
de sa \omterm{def:b2:reaction-free-energy:inversion-temperature}{température d'inversion}, $\boldsymbol{T_i \approx \qty{1.11e3}{K}}$
(environ \qty{840}{\celsius}).
\end{solution}
